\section{主理想整环上的有限生成模}\label{sec:galois-1-10}
\subsection{定义}\label{sec:galois-1-10-1}
\begin{definition*}
设 $M$ 是整环 $A$ 上的模. 若 $\operatorname{Ann}_A(m)\ne\{0\}$，称 $m$ 为\newterm{galois-torsion-element}{挠元}{torsion element}. 挠元组成的子集
\[M_{\mathrm{tors}}:=\{m\in M:\operatorname{Ann}_A(m)\ne\{0\}\}\]
是 $M$ 的 $A$-子模，称为 $M$ 的\newterm{galois-torsion-submodule}{挠子模}{torsion submodule}. 若 $M_{\mathrm{tors}}=\{0\}$，称 $M$ \newterm{galois-torsion-free-module}{无挠}{torsion-free module}；若 $M=M_{\mathrm{tors}}$，称 $M$ 为\newterm{galois-torsion-module}{挠模}{torsion module}.
\end{definition*}
\subsection{主理想整环上的挠模：准素分解}\label{sec:galois-1-10-2}
现设 $A$ 是主理想整环. 回顾，在这种情形下，每个非零素理想都是极大理想，即 $\operatorname{Max}(A)=\operatorname{Spec}(A)\setminus\{0\}$. 这样的理想 $\mathfrak p$ 由一个不可约元 $p$ 生成，而 $p$-进赋值仅依赖于 $\mathfrak p$. 将其记作 $\nu_{\mathfrak p}\colon A\setminus\{0\}\longrightarrow\mathbb N$. 此外，对 $A$ 的每个非零理想 $I$，它的一个生成元 $x$ 的赋值 $\nu_{\mathfrak p}(x)$ 不依赖于生成元的选择，将它记为 $\nu_{\mathfrak p}(I)$. 用理想表示，$x$ 的分解就写成
\[I=\prod_{\mathfrak p}\mathfrak p^{\nu_{\mathfrak p}(I)}.\]
\begin{proposition*}
设 $M$ 是主理想整环 $A$ 上的有限生成挠模，其零化理想为 $I:=\operatorname{Ann}_A(M)$. 则
\begin{enumerate}[label=\roman*)]
\item 对每个 $\mathfrak p\in\operatorname{Max}(A)$，有
\[M[\mathfrak p^\infty]=M[\mathfrak p^{\nu_{\mathfrak p}(I)}]\simeq M/\mathfrak p^{\nu_{\mathfrak p}(I)}\simeq M_{\mathfrak p}.\]
\item 有 $M=\bigoplus_{\mathfrak p\in\operatorname{Max}(A)}M[\mathfrak p^\infty]$，其中只有有限个非零项.
\item 存在 $m\in M$ 使 $\operatorname{Ann}_A(m)=\operatorname{Ann}_A(M)$.
\end{enumerate}
\end{proposition*}
\noindent\textit{译注：} i) 中的商模原文写作 $M/\mathfrak p^{\nu_{\mathfrak p}(I)}$，未写分母中的 $M$，此处保留原式.
\begin{proof}
先证明 $I\ne0$. 取 $M$ 的一个生成族 $m_1,\ldots,m_s$，则 $I=\bigcap_{i=1}^s\operatorname{Ann}_A(m_i)$. 由假设，每个 $\operatorname{Ann}_A(m_i)$ 都非零；由于 $A$ 是整环，它们的乘积非零，更不用说它们的交非零.

于是等式 $I=\prod_{\mathfrak p}\mathfrak p^{\nu_{\mathfrak p}(I)}$ 表明 $I$ 是有限个极大理想的乘积. 因此可以应用第~\ref{sec:galois-1-9-2}~分目的定理及其后的注，得到 i) 和 ii).

对于 iii)，先注意，若 $m=\sum_{\mathfrak p}m_{\mathfrak p}$，其中 $m_{\mathfrak p}\in M[\mathfrak p^\infty]$ 非零，则
\[\operatorname{Ann}_A(m)=\bigcap_{\mathfrak p}\operatorname{Ann}_A(m_{\mathfrak p}),\]
因为各个 $M[\mathfrak p^\infty]$ 构成直和. 此外，各个 $\operatorname{Ann}_A(m_{\mathfrak p})$ 两两互素，故
\[\bigcap_{\mathfrak p}\operatorname{Ann}_A(m_{\mathfrak p})=\prod_{\mathfrak p}\operatorname{Ann}_A(m_{\mathfrak p}).\]
同样有 $\operatorname{Ann}_A(M)=\prod_{\mathfrak p}\operatorname{Ann}_A(M[\mathfrak p^\infty])$，这说明
\[\operatorname{Ann}_A(M[\mathfrak p^\infty])=\mathfrak p^{\nu_{\mathfrak p}(I)},\]
因此 $M[\mathfrak p^{\nu_{\mathfrak p}(I)}]\setminus M[\mathfrak p^{\nu_{\mathfrak p}(I)-1}]\ne\varnothing$. 在后一个集合中取 $m_{\mathfrak p}$，便有 $\operatorname{Ann}_A(m_{\mathfrak p})=\mathfrak p^{\nu_{\mathfrak p}(I)}$，进而 $\operatorname{Ann}_A(m)=I=\operatorname{Ann}_A(M)$.
\end{proof}
\subsection{主理想整环上的无挠模}\label{sec:galois-1-10-3}
先设 $A$ 是分式域为 $K:=\operatorname{Frac}(A)$ 的整环，以 $\iota_M$ 表示将 $m$ 映到 $1\otimes m$ 的 $A$-线性映射 $\iota_M\colon M\longrightarrow K\otimes_A M$.
\begin{lemma*}
$M_{\mathrm{tors}}=\operatorname{Ker}(\iota_M)$，而 $M/M_{\mathrm{tors}}\simeq\operatorname{Im}(\iota_M)$ 无挠，并生成 $K$-向量空间 $K\otimes_A M$.
\end{lemma*}
\begin{proof}
以 $S$ 表示乘法子集 $A\setminus\{0\}$. 我们已知 $K\otimes_A M$ 可以与 $S^{-1}M=\{m/s:m\in M,\ s\in S\}$ 等同，且 $\iota_M(m)=m/1$. 由 $S^{-1}M$ 的定义，$m/1=0$ 当且仅当存在 $t\in S$ 使 $tm=0$，这等价于 $m\in M_{\mathrm{tors}}$.

与任意同态一样，$\iota_M$ 诱导 $M/\operatorname{Ker}(\iota_M)$ 到 $\operatorname{Im}(\iota_M)$ 的同构. 这就将 $M/M_{\mathrm{tors}}$ 与一个 $K$-向量空间的 $A$-子模等同，而这样的子模显然无挠. 最后，由构造，$\iota_M(M)$ 在 $K$ 上生成 $K\otimes_A M$.
\end{proof}
因此，当 $M$ 无挠时，它通过 $\iota_M$ 与 $K$-向量空间 $M_K:=K\otimes_A M$ 的一个 $A$-子模等同. 我们将略去记号中的 $\iota_M$，滥用记号写成 $M\subset M_K$. 可以将 $M_K$ 看作``由 $M$ 生成的 $K$-向量空间''. 更确切地说：
\begin{lemma*}
设 $V$ 是 $K$-向量空间，$M$ 是 $V$ 的一个 $A$-子模. 则 $M$ 无挠，而典范线性映射 $M_K\longrightarrow V$、$k\otimes m\mapsto km$ 给出 $M_K$ 到 $V$ 中由 $M$ 生成的 $K$-子空间
\[KM=\{v\in V:\exists s\in A\setminus\{0\},\ sv\in M\}\]
的同构.
\end{lemma*}
\begin{proof}
$M$ 显然无挠. 回顾，所说的映射确实存在，因为它是由标量扩张的泛性质对应于包含映射 $M\subset V$ 的映射. 注意，$M_K=S^{-1}M$ 的每个元素都形如 $\frac1s\otimes m=m/s$，其中 $s\in A\setminus\{0\}$、$m\in M$. 它在 $V$ 中的像为 $\frac1s\cdot m$，只有 $m=0$ 时才为零. 因此 $M_K\longrightarrow V$ 是单射. 另一方面，其像显然是由 $M$ 生成的 $K$-子空间 $KM$，并具有所述形式.
\end{proof}
对于有限生成 $A$-模 $M$，$K$-向量空间 $M_K:=K\otimes_A M$ 也有限生成，故有限维.
\begin{theorem*}
设 $A$ 是主理想整环，$M$ 是有限生成无挠模. 则 $M$ 是自由 $A$-模，秩为 $\dim_K(M_K)$.
\end{theorem*}
\begin{example*}
为说明``主理想整环''这一假设的重要性，考虑 $A=\mathbb C[X,Y]$ 时的反例. 我们已知理想 $M=(X,Y)\subset A$ 不是自由 $A$-模，但它无挠. 它的生成族至少有两个元素，而 $M_K=K$ 的维数为 $1$.
\end{example*}
\begin{remark*}
事实上，若整环 $A$ 满足上述定理的结论，则它是主理想整环（考虑定理中 $n=1$ 的情形）.
\end{remark*}
\begin{corollary*}
在主理想整环 $A$ 上，每个有限生成模 $M$ 都有分解 $M=M_{\mathrm{tors}}\oplus M'$，其中 $M'$ 是秩为 $\dim_K(K\otimes_A M)$ 的自由模.
\end{corollary*}
\begin{proof}
上述定理和引理表明，$M/M_{\mathrm{tors}}$ 是秩为 $n:=\dim_K(K\otimes_A M)$ 的自由模. 取 $M$ 中的元素 $e_1,\ldots,e_n$，使其像 $\overline e_1,\ldots,\overline e_n$ 构成 $M/M_{\mathrm{tors}}$ 的一组基. 这个族在 $A$ 上线性无关（因为在商中线性无关），故生成一个秩为 $n$ 的自由子模 $M'$，且投影 $\pi\colon M\longrightarrow M/M_{\mathrm{tors}}$ 在 $M'$ 上的限制是同构 $M'\xrightarrow{\sim}M/M_{\mathrm{tors}}$. 特别地，$M_{\mathrm{tors}}\cap M'=\{0\}$. 此外，若 $m\in M$ 的像为 $\overline m=\sum_i a_i\overline e_i$，令 $m'=\sum_i a_ie_i$，则 $m-m'\in\operatorname{Ker}(\pi)=M_{\mathrm{tors}}$，故 $M=M_{\mathrm{tors}}+M'$.
\end{proof}
根据上面的第二个引理，上述定理等价于下面这个更强定理的 i).
\subsection{定理}\label{sec:galois-1-10-4}
\begin{theorem*}
设 $A$ 是主理想整环，$V$ 是 $n$ 维 $K$-向量空间.
\begin{enumerate}[label=\roman*)]
\item $V$ 的每个有限生成且在 $K$ 上生成 $V$ 的 $A$-子模 $M$ 都是秩为 $n$ 的自由模.
\item 若 $M'\subset M$ 是在 $K$ 上生成 $V$ 的 $A$-子模，则
\begin{enumerate}[label=(\alph*)]
\item 存在 $M$ 的一组基 $e_1,\ldots,e_n$ 及 $A$ 中元素 $a_1\mid a_2\mid\cdots\mid a_n$，使 $a_1e_1,\ldots,a_ne_n$ 是 $M'$ 的一组基.
\item 理想 $(a_1)\supseteq\cdots\supseteq(a_n)$ 仅依赖于 $M$ 与 $M'$. 更确切地说，对从 $M$ 的一组基到 $M'$ 的一组基的任意过渡矩阵 $P$，以及任意 $1\le r\le n$，有
\[(a_1a_2\cdots a_r)=\text{由 $P$ 的全部 $r\times r$ 子式生成的理想}.\]
\end{enumerate}
\end{enumerate}
\end{theorem*}
\begin{proof}
对 $n$ 归纳.

设 $n=1$，固定非零 $v\in V$. 取 $M$ 的生成元 $m_1,\ldots,m_r$. 因 $V$ 是一维的，可写成 $m_i=(a_i/s_i)v$. 令 $w:=\frac1{s_1\cdots s_r}v\ne0$. 对每个 $i$ 都有 $m_i\in Aw\subset V$，故 $M\subset Aw$. 映射 $\Psi\colon A\longrightarrow Aw$、$a\mapsto aw$ 是 $A$-模同构. 因此 $\Psi^{-1}(M)$ 是 $A$ 的理想，故为主理想，由某个 $x\in A$ 生成. 因 $A$ 是整环，$\Psi^{-1}(M)$ 实际上是以 $\{x\}$ 为基的自由 $A$-模，从而 $M$ 是以 $e_1:=xw$ 为基的自由 $A$-模. 同样，$\Psi^{-1}(M')$ 是包含 $(x)$ 的理想，故形如 $(a_1x)$，从而 $M'$ 是以 $a_1xw=a_1e_1$ 为基的自由 $A$-模.\footnote{译注：原文写作``包含 $(x)$''，与 $M'\subset M$ 所给的包含方向相反，此处保留原文.}

现设 $n>1$，取 $V$ 中的一条 $K$-直线 $D$. 令 $L:=M\cap D$. 它是 $M$ 的 $A$-子模，因 $A$ 为 Noether 环而有限生成；它生成直线 $D$，因为对每个 $v\in D$，存在 $s\in A\setminus\{0\}$ 使 $sv\in M$，从而 $sv\in L=M\cap D$，且 $v=\frac1s sv\in KL$. 同样，$L':=M'\cap D$ 是 $L$ 的有限生成 $A$-子模，且在 $K$ 上生成 $D$. 刚才已证明，这些模是秩 $1$ 的自由 $A$-模，并且可以选取 $L$ 的一个基 $e_1$ 及 $a_1\in A$，使 $a_1e_1$ 是 $L'$ 的基.

现令 $N$ 为 $M$ 在商空间 $W:=V/D$ 中的像. 它是 $W$ 的有限生成 $A$-子模，且在 $K$ 上生成 $W$，因为它是一个生成空间的模的像. $M'$ 在 $W$ 中的像 $N'$ 也具有这一性质. 因此根据归纳假设，可以设存在 $N$ 的 $A$-基 $\overline e_2,\ldots,\overline e_n$ 及元素 $a_2\mid a_3\mid\cdots\mid a_n$，使 $a_2\overline e_2,\ldots,a_n\overline e_n$ 是 $N'$ 的 $A$-基.

取 $M$ 中的元素 $e_2,\ldots,e_n$，使它们在 $N$ 中的像分别为 $\overline e_2,\ldots,\overline e_n$. 则 $(e_1,\ldots,e_n)$ 是 $M$ 的基. 事实上，这个族是 $V$ 的 $K$-基，因而更是 $A$-线性无关的. 此外，若 $m\in M$ 在 $N$ 中的像为 $\overline m=\sum_{i=2}^n b_i\overline e_i$，则 $m-\sum_{i=2}^n b_ie_i$ 属于投影 $V\longrightarrow W=V/D$ 的核 $D$，从而属于 $D\cap M=L$，故形如 $b_1e_1$. 这说明 $(e_1,\ldots,e_n)$ 也生成 $M$，因此是它的 $A$-基. i) 得证.

证明 ii)(a) 以前，注意一般没有理由认为元素 $a_ie_i$（$i=2,\ldots,n$）属于 $M'$. 对每个 $i$，起初只能断言 $a_ie_i\in M'+D$，因而也有
\[a_ie_i\in M'+\bigl(D\cap(M'+a_iM)\bigr)=M'+a_iL.\]
不过，假设可以选取直线 $D$ 使 $a_1\mid a_2$. 这时 $L'=a_1L\supset a_iL$，故对每个 $i$ 有 $a_ie_i\in M'$，与上面一样，族 $(a_1e_1,\ldots,a_ne_n)$ 是 $M'$ 的基.

\enlargethispage{2pt}
因此，为证明 ii)(a)，只需证明可以选取直线 $D$ 使 $a_1\mid a_2$. 为此，取 $t\in A$ 使 $tM\subset M'$（这样的 $t$ 存在，因为 $M'$ 生成 $V$ 且 $M$ 有限生成）. 商模 $M'/tM$ 是挠模，因为它被 $t$ 零化. 以 $t'$ 表示其零化理想的一个生成元：
\[(t')=\operatorname{Ann}_A(M'/tM)=\{a\in A:aM'\subset tM\}.\]
注意 $t'\mid t$. 根据第~\ref{sec:galois-1-10-2}~分目命题的 iii)，存在 $m'\in M'$ 使
\[(t')=\operatorname{Ann}_A(m'+tM)=\{a\in A:am'\in tM\}.\]
取 $D$ 为 $m'$ 生成的直线. 有 $tL=t(D\cap M)=D\cap tM$，因此也有 $tL=L'\cap tM$. 从而对每个 $l'\in L'$ 及 $a\in A$，有 $al'\in tL\Longleftrightarrow al'\in tM$. 由此推出
\[(t')=\{a\in A:am'\in tL\}=\{a\in A:aL'\subset tL\}.\]
取 $L$ 的基 $e_1$ 及 $a_1$，使 $a_1e_1$ 是 $L'$ 的基. 以 $\Psi$ 表示同构 $A\xrightarrow{\sim}L$、$a\mapsto ae_1$，则 $\Psi^{-1}(L')=(a_1)$. 还有 $\Psi^{-1}(tL)=(t)$，而等式
\[(t')=\{a\in A:a(a_1)=(t)\}\]
蕴含 $(t')(a_1)=(t'a_1)=(t)$.\footnote{译注：集合条件中的等号为原文所印；与前一式对照，此处的理想条件疑应为包含关系.}
现以 $t''$ 表示 $N'/tN$ 的零化理想的一个生成元：
\[(t'')=\operatorname{Ann}_A(N'/tN)=\{a\in A:aN'\subset tN\}.\]
因为 $N'/tN$ 是 $M'/tM$ 的商，有 $t''\mid t'$. $N$ 的基 $\overline e_2,\ldots,\overline e_n$ 给出同构 $\Phi\colon A\oplus\cdots\oplus A\xrightarrow{\sim}N$，使
\[\Phi^{-1}(N')=(a_2)\oplus\cdots\oplus(a_n),\qquad \Phi^{-1}(tN)=(t)\oplus\cdots\oplus(t).\]
由此可知 $a_i$ 整除 $t$，且 $(t'')=\bigcap_{i=2}^n(t_i'')$，其中 $t=a_it_i''$. 特别地，$t_2''$ 整除 $t''$，从而整除 $t'$. 于是等式 $(a_1t')=(a_2t_2'')=(t)$ 保证 $a_1$ 整除 $a_2$，正如所需.

现在证明 ii)(b). 对矩阵 $P\in\mathrm M_n(A)$，以 $I_r(P)$ 表示它的所有 $r\times r$ 子式在 $A$ 中生成的理想. 设 $P$ 是从 $M$ 的基 $\varepsilon_1,\ldots,\varepsilon_n$ 到 $M'$ 的基 $\varepsilon_1',\ldots,\varepsilon_n'$ 的过渡矩阵. 因此 $P\in\mathrm M_n(A)\cap\mathrm{GL}_n(K)$. 以 $Q$ 表示从 $e_1,\ldots,e_n$ 到 $\varepsilon_1,\ldots,\varepsilon_n$ 的过渡矩阵，以 $Q'$ 表示从 $\varepsilon_1',\ldots,\varepsilon_n'$ 到 $a_1e_1,\ldots,a_ne_n$ 的过渡矩阵. 则 $Q,Q'\in\mathrm{GL}_n(A)$，且
\[QPQ'=\Delta:=\operatorname{diag}(a_1,\ldots,a_n),\]
它是从 $e_i$ 到 $a_ie_i$ 的过渡矩阵. 特别地，$\Delta$ 的各项都是 $P$ 各项的 $A$-线性组合，因此 $I_1(\Delta)\subset I_1(P)$. 但也可以写成 $\Delta=Q^{-1}PQ'^{-1}$；\footnote{译注：原文在此仍写 $\Delta=Q^{-1}PQ'^{-1}$，而由前式应解出 $P=Q^{-1}\Delta Q'^{-1}$；保留原式.}因为 $Q^{-1}$、$Q'^{-1}$ 的系数在 $A$ 中，又有 $I_1(P)\subset I_1(\Delta)$. 因 $I_1(\Delta)=(a_1)$，得到 $I_1(P)=(a_1)$.

当 $r>1$ 时可作类似论证. 对方阵 $R=(a_{ij})\in\mathrm M_n(A)$ 以及 $\{1,\ldots,r\}$ 中两个基数为 $r$ 的子集 $I,J$，以 $a_{IJ}$ 表示对应于 $I,J$ 的子式，即矩阵 $R_{IJ}\in\mathrm M_r(A)$ 的行列式；后者的行（相应地，列）取自编号属于 $I$（相应地，$J$）的行（相应地，列）.\footnote{译注：原文的下标集合写作 $\{1,\ldots,r\}$；与 $n$ 阶矩阵的全部 $r$ 阶子式对照，疑应为 $\{1,\ldots,n\}$.}各个 $a_{IJ}$ 构成一个 $\binom nr$ 阶矩阵的各项，将它记作 $\bigwedge^r R$（须为 $I,J$ 选定一个次序，例如字典序）. 根据定义，$I_r(R)=I_1(\bigwedge^r R)$. 现在有等式
\[(\bigwedge^r Q)(\bigwedge^r P)(\bigwedge^r Q')=\bigwedge^r\Delta\]
（注意，当 $r=n$ 时，这仅表示 $\det(Q)\det(P)\det(Q')=\det(\Delta)$）. 其中 $\bigwedge^rQ\in\mathrm{GL}_{\binom nr}(A)$，其逆为 $\bigwedge(Q^{-1})$，$Q'$ 也同样如此. 因而
\[I_r(P)=I_1(\bigwedge^rP)=I_1(\bigwedge^r\Delta)=I_r(\Delta).\]
而立刻可见 $I_r(\Delta)=(a_1\cdots a_r)$.
\end{proof}
\begin{remark*}
也可以从 $M$ 的一组基到 $M'$ 的一组基的过渡矩阵 $P$ 出发，调整 Gauss 消元法，利用行变换（相应地，列变换）将其化为对角矩阵 $\operatorname{diag}(a_1,\ldots,a_n)$，从而证明定理 ii)(a) 中的基 $e_1,\ldots,e_n$ 存在；这些变换由左乘（相应地，右乘）$\mathrm{GL}_n(A)$ 中的矩阵给出.
\end{remark*}

\subsection{挠模：不变因子}\label{sec:galois-1-10-5}
若模 $M$ 形如 $A/I$，其中 $I$ 是非零理想，就称它是\newterm{galois-cyclic-module}{循环模}{cyclic module}.
\begin{corollary*}
设 $M$ 是主理想整环 $A$ 上的有限生成挠模. 则 $M$ 是循环模的直和. 更确切地说，存在唯一的整数 $r$ 及唯一的非零真理想列 $I_1\supset I_2\supset\cdots\supset I_r$，使
\[M\simeq A/I_1\oplus A/I_2\oplus\cdots\oplus A/I_r.\]
理想 $I_i$ 称为 $M$ 的\newterm{galois-invariant-factor}{不变因子}{invariant factor}.
\end{corollary*}
\begin{proof}
先证存在性. 取生成元 $m_1,\ldots,m_r$，使 $r$ 最小. 考虑将 $(a_1,\ldots,a_r)$ 映到 $\sum_i a_im_i$ 的 $A$-模同态 $A^r\longrightarrow M$. 它是满射，故以 $K$ 表示其核，它就诱导同构 $A^r/K\xrightarrow{\sim}M$. 根据上个定理，存在 $A^r$ 的一组基 $e_1,\ldots,e_r$（不一定是标准基）以及元素 $a_1\mid\cdots\mid a_r$，使 $a_1e_1,\ldots,a_re_r$ 是 $K$ 的一组基. 因此 $A^r/K\simeq\bigoplus_{i=1}^r A/(a_i)$. 存在性得证.

现证唯一性，即从 $M$ 恢复 $r$ 和 $I_i$. 对每个 $\mathfrak p\in\operatorname{Max}(A)$，考虑递减的 $A$-子模列
\[(\mathfrak p^nM)_{n\in\mathbb N}=M\supset\mathfrak pM\supset\mathfrak p^2M\supset\cdots.\]
相继商 $\mathfrak p^nM/\mathfrak p^{n+1}M$ 是被 $\mathfrak p$ 零化的 $A$-模，因而是域 $A/\mathfrak p$ 上的向量空间. 对有限生成的 $M$，维数
\[d_{n,\mathfrak p}(M):=\dim_{A/\mathfrak p}(\mathfrak p^nM/\mathfrak p^{n+1}M)\]
有限，并且当 $n\gg0$ 时等于零. 若 $M$ 是 $\mathfrak q$-挠模，其中 $\mathfrak q\ne\mathfrak p$，则对每个 $n$ 有 $\mathfrak p^nM=M$，故 $d_{n,\mathfrak p}(M)=0$. 因为 $d_{n,\mathfrak p}(M\oplus M')=d_{n,\mathfrak p}(M)+d_{n,\mathfrak p}(M')$，一般地可得
\[d_{n,\mathfrak p}(M)=d_{n,\mathfrak p}(M[\mathfrak p^\infty]).\]
特别地，
\[d_{n,\mathfrak p}(A/I)=d_{n,\mathfrak p}(A/\mathfrak p^{\nu_{\mathfrak p}(I)}).\]
现对 $m\in\mathbb N$ 计算得到
\begin{align*}
\mathfrak p^n(A/\mathfrak p^m)&\simeq\begin{cases}\mathfrak p^n/\mathfrak p^m,&n\le m,\\0,&n\ge m,\end{cases}\\
\mathfrak p^n(A/\mathfrak p^m)/\mathfrak p^{n+1}(A/\mathfrak p^m)&\simeq\begin{cases}\mathfrak p^n/\mathfrak p^{n+1},&n<m,\\0,&n\ge m.\end{cases}
\end{align*}
而 $\mathfrak p^n/\mathfrak p^{n+1}$ 是一维 $A/\mathfrak p$-向量空间. 事实上，若 $p$ 是 $\mathfrak p$ 的一个生成元，则 $p^n$ 是 $\mathfrak p^n$ 的一个生成元，所以它的像生成 $A$-模 $\mathfrak p^n/\mathfrak p^{n+1}$，因而也生成这个 $A/\mathfrak p$-向量空间. 因此
\[d_{n,\mathfrak p}(A/\mathfrak p^m)=\begin{cases}1,&n<m,\\0,&n\ge m.\end{cases}\]
从而，对形如 $M=A/\mathfrak p^{m_1}\oplus A/\mathfrak p^{m_2}\oplus\cdots\oplus A/\mathfrak p^{m_s}$ 的模，其中整数 $m_1\ge\cdots\ge m_s$，序列 $(d_{n,\mathfrak p}(M))_{n\in\mathbb N}$ 递减，且可以通过不等式
\[d_{m_i,\mathfrak p}(M)<i\le d_{m_i-1,\mathfrak p}(M)\]
恢复各个 $m_i$. 这些不等式只涉及序列 $(d_{n,\mathfrak p}(M))_{n\in\mathbb N}$，并说明只要 $s>d_{0,\mathfrak p}(M)$，就有 $m_s=0$.\footnote{译注：原文在此沿用上述不等式讨论零指数，但此时下标 $m_s-1=-1$ 不在所定义的范围内；保留原论述.}

特别地，对于模 $M=A/I_1\oplus\cdots\oplus A/I_r$，其中 $I_1\subset\cdots\subset I_r\ne A$，可以由公式
\[r:=\operatorname{Max}\{d_{0,\mathfrak p}(M):\mathfrak p\in\operatorname{Max}(A)\}\]
恢复 $r$，再由以下公式恢复 $I_i$：
\[I_i=\prod_{\mathfrak p\in\operatorname{Max}(A)}\mathfrak p^{m_i},\qquad
\text{其中 $m_i$ 由 }d_{m_i,\mathfrak p}(M)<i\le d_{m_i-1,\mathfrak p}(M)\text{ 确定}.\]
将 $I_i$ 重新编号为 $I_{r-i}$（用递减列代替递增列），即可得到定理中的唯一性.\footnote{译注：原文重编号为 $I_{r-i}$，下标在 $i=r$ 时成为 $0$；这里保留原下标.}
\end{proof}
对于 $A$ 的一个极大理想 $\mathfrak p$，单独强调 $\mathfrak p$-挠模这一特例也很有用.
\begin{corollary*}
设 $A$ 是主理想整环，$M$ 是有限生成的 $\mathfrak p$-挠 $A$-模. 存在唯一的整数 $r$ 及唯一的序列 $m_1\le m_2\le\cdots\le m_r$，使
\[M\simeq A/\mathfrak p^{m_1}\oplus A/\mathfrak p^{m_2}\oplus\cdots\oplus A/\mathfrak p^{m_r}.\]
\end{corollary*}
\subsection{在线性自同态约化中的应用}\label{sec:galois-1-10-6}
设 $K$ 是域，$V$ 是 $n$ 维向量空间，$u\in\operatorname{End}_K(V)$ 是 $V$ 的线性自同态. 将 $X$ 映到 $u$ 的唯一 $K$-代数同态 $K[X]\longrightarrow\operatorname{End}_K(V)$ 赋予 $V$ 一个 $K[X]$-模结构. 这个同态的核 $I$ 是 $K[X]$ 的理想，因此由唯一的首一多项式 $f_u\in K[X]$ 生成，称为 $u$ 的\newterm{galois-minimal-polynomial-endomorphism}{极小多项式}{minimal polynomial of an endomorphism}. 设 $f_u=f_1^{v_1}\cdots f_r^{v_r}$ 是 $f_u$ 分解为首一不可约多项式幂的乘积. $V$ 作为 $K[X]$-模的准素分解
\[V=V[f_1^{v_1}]\oplus\cdots\oplus V[f_r^{v_r}]\]
正是由核分解引理给出的分解，因为根据定义，
\[V[f_i^{v_i}]=\operatorname{Ker}(f_i(u)^{v_i}).\]
假设 $K$ 代数闭，因而 $f_i=X-\lambda_i$，其中 $\lambda_i\in K$ 两两不同. 则 $V[f_i^{v_i}]$ 是 $u$ 对应于特征值 $\lambda_i$ 的\newterm{galois-generalized-eigenspace}{广义特征子空间}{generalized eigenspace}，而 $V[f_i]$ 是其特征子空间. 现在将循环子模的直和分解写成
\begin{align*}
V[f_i^{v_i}]&=V_{i1}\oplus\cdots\oplus V_{ir_i},\\
V_{ij}&=K[X]w_{ij},\qquad K[X]/(f_i^{m_{ij}})\xrightarrow{\sim}K[X]w_{ij}.
\end{align*}
则每个 $V_{ij}$ 都在 $u$ 下稳定，且族
\[w_{ijk}:=(u-\lambda_i)^{k-1}(w_{ij}),\qquad k=1,\ldots,m_{ij}\]
是 $V_{ij}$ 的一组 $K$-基，$u$ 在这组基下的矩阵形如 $\lambda_iI+N$，其中 $N$ 是 $m_{ij}\times m_{ij}$ 矩阵，其``上副对角线''上的元素为 $1$，其余位置为 $0$.\footnote{译注：原文如此指定基的排列及上副对角线；在通常的列向量约定下，按这里的 $k$ 递增排列得到的 $1$ 位于下副对角线，反向排列才得到上副对角线. 此处保留原文.}换句话说，$u$ 在基
\[\{w_{ijk}:i=1,\ldots,r,\ j=1,\ldots,r_i,\ k=1,\ldots,m_{ij}\}\]
下的矩阵具有 Jordan 形.

